% Préambule commun aux cours IPR Scholar.
\documentclass[11pt,a4paper,oneside]{article}
\usepackage[T1]{fontenc}
\usepackage[utf8]{inputenc}
\usepackage[french]{babel}
\usepackage{lmodern}
\usepackage{microtype}
\usepackage[a4paper,margin=2.35cm]{geometry}
\usepackage{amsmath,amssymb,mathtools}
\usepackage{booktabs,longtable,array}
\usepackage{xcolor}
\usepackage{hyperref}
\usepackage{tikz}
\usetikzlibrary{arrows.meta,positioning,shapes.geometric,calc}
\hypersetup{colorlinks=true,linkcolor=blue!55!black,citecolor=blue!55!black,urlcolor=blue!55!black}
\definecolor{iprblue}{HTML}{1F4E79}
\definecolor{iprgreen}{HTML}{2D6A4F}
\newcommand{\IPR}{\textsc{ipr}}
\newcommand{\coursmeta}[2]{\author{Projet \IPR{} -- documentation Scholar}\date{#1}\maketitle\begin{center}\small #2\end{center}\vspace{0.5cm}\tableofcontents\clearpage}
\newcommand{\resultat}[1]{\paragraph{À retenir.} #1}

\title{Cours 5 -- Texture, clarté et filtrage gaussien multi-échelle}
\begin{document}
\coursmeta{3 octobre 2026}{Approche de traitement du signal utilisée pour Texture, Clarté et Flou.}

\section{Théorie}
Un flou gaussien est la convolution d'une image $Y$ avec $G_\sigma$. Sa réponse passe-bas isole les structures dont l'échelle est plus grande que $\sigma$. Le résidu $D_\sigma=Y-G_\sigma(Y)$ est un détail local. Une accentuation contrôlée s'écrit
\begin{equation}
Y'=\max(0,Y+aD_\sigma w(Y)).
\end{equation}
Cette famille d'opérations est voisine de l'\emph{unsharp masking}; le choix de $\sigma$ définit ce que l'on appelle ``détail''~\cite{gonzalezWoods,opencvGaussian}.

\section{Application dans \IPR{}}
Texture appelle un flou à $\sigma=1{,}2$ pixels : les fibres, petits contours et micro-variations sont ciblés. Clarté utilise $\sigma=8$ pixels : les volumes et transitions moyennes sont concernés. Pour clarté, un produit de deux \texttt{smoothstep} réduit l'effet sous $Y\approx0{,}16$ et au-dessus de $Y\approx0{,}72$ : cela protège ombres profondes et hautes lumières, zones où le contraste local forme facilement des halos.

Le curseur Flou est volontairement différent : il applique le gaussien directement au RGB avec $\sigma=\texttt{blur}/2$. L'alpha est préservé, car le flouter fabriquerait des bords transparents artificiels.

\section{Exemple}
Sur un portrait, Texture $+20$ révèle les détails d'un vêtement; Clarté $+20$ rend aussi plus présents les contours du visage et l'arrière-plan. Une clarté trop élevée augmente le bruit et les halos; le masque d'extrêmes limite ce défaut mais ne le supprime pas.

\section{Choix des valeurs}
Les sigmas ne découlent pas d'une loi universelle : ce sont des tailles d'échelle choisies pour séparer visuellement micro-texture et contraste de volume à la résolution d'aperçu. Ils doivent évoluer avec une éventuelle normalisation par taille d'image ou ratio d'affichage. La validation doit inclure bruit, peau, feuillage et ciel.

\resultat{Texture et clarté sont le même opérateur de détail à deux échelles; leur différence est une question de fréquence spatiale, pas de nature physique.}
\bibliographystyle{plain}\bibliography{references}
\end{document}
